\subsection{上同伦子的迷向群}

设 G 与 H 为群胚，\(\varphi\) 与 \(\psi\) 为从 H 到 G 的群胚态射。本节的目的，是计算上同伦子 \(\mathrm{Coh}(\varphi, \psi)\) 的迷向群。我们沿用第 3 号中的记号 \(G_1, h_1, \theta_1, \alpha, h\) 。

以 \(\Gamma_{0}\) 表示偶 \((\varphi,\psi)\) 的框架。回顾（II, 第 185 页, 定义 3），这是箭图 \((\mathrm{Orb}(\mathrm{G}),\mathrm{Orb}(\mathrm{H}),\varphi_{0},\psi_{0})\) ，其中 \(\varphi_{0}\) 与 \(\psi_{0}\) 是从 \(\mathrm{Orb}(\mathrm{H})\) 到 \(\mathrm{Orb}(\mathrm{G})\) 的、由映射 \(\varphi\) 与 \(\psi\) 通过取轨道诱导的映射。

在本节的余下部分，我们还将假设箭图 \(\Gamma_{0}\) 连通且非空；据 II, 第 185 页, 命题 4，这等价于假设群胚 \(\mathrm{Coh}(\varphi,\psi)\) 是传递的，或者等价地，箭图 \(G_{1}\) 连通且非空（II, 第 185 页, 命题 2）。

定义 4. —— 我们把下列数据称为偶 \((\varphi, \psi)\) 的基本配备：

\begin{enumerate}
\def\labelenumi{(\roman{enumi})}
\item
  对每个 \(i \in \mathrm{Orb}(\mathbf{G})\) ，取一个顶点 \(a(i)\) ，位于轨道 \(i\) （\(\mathbf{G}\) 的轨道）之中；
\item
  对每个 \(j \in \mathrm{Orb}(\mathrm{H})\) ，取一个顶点 \(b(j)\) ，位于轨道 \(j\) （H 的轨道）之中；
\item
  对每个 \(j \in \text{Orb(H)}\) ，取 G 中的箭 \(c_{1}(j)\) 与 \(c_{2}(j)\) ，分别连接 \(\varphi(b(j))\) 与 \(a(\varphi_{0}(j))\) 、\(\psi(b(j))\) 与 \(a(\psi_{0}(j))\) ；
\item
  框架 \(\Gamma_{0}\) 的一个子箭图 T，其相伴图是图 \(\widetilde{\Gamma}_{0}\) 的极大树；
\item
  G 的一个轨道 \(i_{0}\) 。
\end{enumerate}

选取一个基本配备 \((a,b,c_{1},c_{2},\mathrm{T},i_{0})\) （偶 \((\varphi,\psi)\) 的）。我们定义箭图态射 \(\tau_{1}\) （从 \(\Gamma_{0}\) 到 \(\operatorname{Grp}(\mathrm{G}_{1})\) ）如下：令 \(\tau_{1}(i)=a(i)\) （对 \(i\in\operatorname{Som}(\Gamma_{0})=\operatorname{Orb}(\mathrm{G})\) ），并令

\[
\tau_ {1} (j) = c _ {1} (j) ^ {- 1} \cdot h _ {1} (b (j)) \cdot c _ {2} (j)
\]

（对 \(j \in \mathrm{Fl}(\Gamma_{0}) = \mathrm{Orb}(\mathrm{H})\)）。以 \(\tau_{0}\) 表示 \(\tau_{1}\) 与典范态射 \(\theta_{1}\) （从 \(\mathrm{Grp}(\mathrm{G}_{1})\) 到 \(\mathrm{Coh}(\varphi, \psi)\) ）的复合；它是从 \(\Gamma_{0}\) 到 \(\mathrm{Coh}(\varphi, \psi)\) 的箭图态射。

对 \(i \in \mathrm{Orb}(\mathrm{G})\) ，以 \(\alpha_i \colon \mathrm{G}_{a(i)} \to \mathrm{Coh}(\varphi, \psi)_{a(i)}\) 表示由态射 \(\alpha \colon \mathrm{G} \to \mathrm{Coh}(\varphi, \psi)\) 通过限制到 \(a(i)\) 处的迷向群而得的群同态。

对 \(j\in \mathrm{Orb(H)}\) ，以

\[
\varphi_ {j} = \mathrm{Int} (c _ {1} (j)) ^ {- 1} \circ \varphi_ {b (j)} \colon \mathrm{H} _ {b (j)} \to \mathrm{G} _ {a (\varphi_ {0} (j))}
\]

及

\[
\psi_ {j} = \mathrm{Int} (c _ {2} (j)) ^ {- 1} \circ \psi_ {b (j)} \colon \mathrm{H} _ {b (j)} \to \mathrm{G} _ {a (\psi_ {0} (j))},
\]

表示相应的同态，使得对每个元素 \(f\) ∈ \(\mathrm{H}_{b(j)}\) ，有

\[
\varphi_ {j} (f) = c _ {1} (j) ^ {- 1} \varphi (f) c _ {1} (j) \quad \text { and } \quad \psi_ {j} (f) = c _ {2} (j) ^ {- 1} \psi (f) c _ {2} (j).\tag{6}
\]

对 \(\Gamma_{0}\) 的每个顶点 i，仍以 \(d_{i}\) 表示连接 \(i_{0}\) 与 i 的、树 \(\widetilde{T}\) 中的唯一道路类；把它看作 \(\mathrm{Grp}(\Gamma_{0})\) 中的一条箭。再以 \(\delta_{i}\) 表示 \(\mathrm{Coh}(\varphi,\psi)\) 中的这样一条箭：它是 \(d_{i}\) 在从 \(\mathrm{Grp}(\Gamma_{0})\) 到 \(\mathrm{Coh}(\varphi,\psi)\) 的、由 \(\tau_{0}\) 导出的典范态射下的像；\(\delta_{i}\) 的起点是 \(a(i_{0})\) ，其靶是 \(a(i)\) 。

箭图态射 \(\tau_0\) 、群同态 \(\alpha_i\) （对 \(i \in \mathrm{Orb}(\mathbf{G})\) ）、\(\varphi_j\) 与 \(\psi_j\) （对 \(j \in \mathrm{Orb}(\mathbf{H})\) ）、以及箭 \(\delta_i\) （在 \(\mathrm{Coh}(\varphi, \psi)\) 中）（对 \(i \in \mathrm{Orb}(\mathbf{G})\) ），称为由基本配备导出。

若 \((\mathrm{G}_{i})_{i\in\mathrm{I}}\) 是一族群，我们以 \(*G_{i}\) 表示它们的自由积；元素 \(g\in G_{i}\) 在从 \(G_{i}\) 到 \(*G_{i}\) 的典范映射下的像记作 {[}g{]} ，在无混淆可能时甚至记作 g。若 S 为集合，我们以 F(S) 表示在 S 上构造的自由群（A, I, 第 84 页）。

命题 6. —— 存在唯一的群同态

\[
\Lambda \colon \big (\underset {i \in \operatorname{Orb} (\mathrm{G})} {*} \mathrm{G} _ {a (i)} \big) * \mathrm{F} (\operatorname{Orb} (\mathrm{H})) \to \operatorname{Coh} (\varphi , \psi) _ {a (i _ {0})}
\]

使得

\begin{enumerate}
\def\labelenumi{(\arabic{enumi})}
\setcounter{enumi}{6}
\tightlist
\item
\end{enumerate}

\[
\Lambda (f) = \delta_ {i} \alpha_ {i} (f) \delta_ {i} ^ {- 1} \quad \text {   for   } i \in \operatorname{Orb} (\mathrm{G}) \text {   and   } f \in \mathrm{G} _ {a (i)},\tag{8}
\]

\[
\Lambda (j) = \delta_ {\varphi_ {0} (j)} \tau_ {0} (j) \delta_ {\psi_ {0} (j)} ^ {- 1} \qquad \text {for } j \in \mathrm{Orb} (\mathrm{H}).
\]

同态 \(\Lambda\) 是满射的；其核是 \((\ast_{i}\mathrm{G}_{a(i)})\ast\mathrm{F}(\mathrm{Orb}(\mathrm{H}))\) 的最小正规子群，它包含元素 \(j\) （属于 \(\mathrm{Fl}(\mathrm{T})\) ）以及诸元素 \(\varphi_{j}(f)j\psi_{j}(f)^{-1}j^{-1}\) （对 \(j\in\mathrm{Orb}(\mathrm{H})\) 及 \(f\in\mathrm{H}_{b(j)}\) ）。

同态 \(\Lambda\) 的存在性与唯一性由自由积与自由群的泛性质得出（A, I, 第 85 页, 命题 8）。

设 A 为诸 \(a(i)\) （对 \(i \in \mathrm{Orb}(\mathrm{G})\) ）的集合，并设 \(\mathrm{G}_{\mathrm{A}}\) 为 G 的以 A 为顶点集的全子群胚。对每个 \(x \in \text{Som(G)}\) ，以 \(\overline{x}\) 表示 \(x\) 在 G 中的轨道，并选取 G 的一条箭 \(d_x\) ，它连接 \(x\) 与 \(a(\overline{x})\) 。偶 \(v\) ——它由映射 \(x \mapsto a(\overline{x})\) （从 \(\text{Som(G)}\) 到 A 中的）以及把 \(f \in \text{Fl}_{x,y}(\text{G})\) 对应到元素 \(d_x^{-1}fd_y\) （属于 \(\text{Fl}_{a(\overline{x}),a(\overline{y})}(\text{G}_\text{A})\) ）的映射组成——是一个群胚态射。由 II, 第 182 页, 命题 1 推出，\(v\) 是同伦等价态射。令 \(\varphi' = v \circ \varphi\) 与 \(\psi' = v \circ \psi\) ，并以 \(w\) 表示从 \(\text{Coh}(\varphi, \psi)\) 到 \(\text{Coh}(\varphi', \psi')\) 的典范态射；它是同伦等价态射（II, 第 187 页, 定理 1）。

\(G_{A}\) 的轨道是诸集合 \(\{a\}\) （对 \(a \in A\) ），且单射 \(G_{A} \to G\) 诱导一个从 \(\mathrm{Orb}(G_{\mathrm{A}})\) 到 \(\mathrm{Orb}(G)\) 上的双射，我们将借此把这两个集合等同起来。定义基本配备 \((a', b', \beta_{1}', \beta_{2}', T', i_{0})\) （偶 \((\varphi', \psi')\) 的）如下：令 \(a'(i) = a(i)\) （对 \(i \in \mathrm{Orb}(G)\) ）、\(b'(j) = b(j)\) 、\(\beta_{1}'(j) = v(c_{1}(j))\) 、\(\beta_{2}'(j) = v(c_{2}(j))\) （对 \(j \in \mathrm{Orb}(H)\) ），并令 \(T' = T\) 。由这个基本配备导出的群同态 \(\varphi_{j}'\) 与 \(\psi_{j}'\) （对 \(j \in \mathrm{Orb}(H)\) ）、箭图态射 \(\tau_{0}'\) 、箭 \(\delta_{i}'\) （对 \(i \in \mathrm{Orb}(G)\) ）、从而还有群同态 \(\Lambda'\) ，分别是相应的同态 \(\varphi_{j}\) 、\(\psi_{j}\) 、箭图态射 \(\tau_{0}\) 、相应的箭 \(\delta_{i}\) 以及同态 \(\Lambda\) 与 w 的复合。

设 B 为诸 \(b(j)\) （对 \(j \in \mathrm{Orb}(\mathrm{H})\) ）的集合，并设 \(\mathrm{H}_{\mathrm{B}}\) 为 H 的以 B 为顶点集的全子群胚；以 \(u: \mathrm{H}_{\mathrm{B}} \to \mathrm{H}\) 表示典范单射；令 \(\varphi'' = \varphi' \circ u\) 与 \(\psi'' = \psi' \circ u\) 。态射 \(u\) 诱导一个双射 \(\mathrm{B} \to \mathrm{Orb}(\mathrm{H})\) ，我们将借此把这两个集合等同起来。我们再由 II, 第 187 页, 定理 1 得到一个典范同伦等价 \(w': \operatorname{Coh}(\varphi'', \psi'') \to \operatorname{Coh}(\varphi', \psi')\)。此外，偶 \((\varphi'', \psi'')\) 配备有一个基本配备 \((a'', b'', \beta_1'', \beta_2'', T'', i_0)\) ，使得 \(\Lambda', \varphi_j', \psi_j', \tau_0', \delta_i' (i, i' \in \operatorname{Orb}(\mathrm{G}), j \in \operatorname{Orb}(\mathrm{H}))\) 分别是 \(w'\) 与 \(\Lambda'', \varphi_j'', \psi_j'', \tau_0'', \delta_i'\) 的复合。

我们用下面的图表总结已引进的各个群胚态射：

\[
\begin{array}{c} \mathrm {H_ {B}} \xrightarrow [ \psi^ {\prime \prime} ]{\varphi^ {\prime \prime}} \mathrm {G_ {A}} \xrightarrow {\alpha^ {\prime \prime}} \mathrm{Coh} (\varphi^ {\prime \prime}, \psi^ {\prime \prime}) \\ \Big \downarrow \qquad \qquad \qquad \Big \| \qquad \qquad \qquad \Big \downarrow w ^ {\prime} \\ \mathrm{H} \xrightarrow [ \psi^ {\prime} ]{\varphi^ {\prime}} \mathrm {G_ {A}} \xrightarrow {\alpha^ {\prime}} \mathrm{Coh} (\varphi^ {\prime \prime}, \psi^ {\prime \prime}) \\ \Big \| \qquad \qquad \qquad \Big | _ {v} \qquad \qquad \qquad \Big | _ {w} \\ \mathrm{H} \xrightarrow [ \psi ]{\varphi} \mathrm{G} \xrightarrow {\alpha} \mathrm{Coh} (\varphi , \psi)  . \end{array}\tag{9}
\]

因此，为证明命题，我们可以假设 \(A = \text{Som}(G)\) 与 \(B = \text{Som}(H)\) ，换言之，典范映射 \(\text{Som}(G) \to \text{Orb}(G)\) 与 \(\text{Som}(H) \to \text{Orb}(H)\) 是双射；我们将在证明的余下部分在此假设之下进行。

这时箭图 \(G_{1}\) 以 A 为其顶点集，以诸集合 \(G_{a}\) （\(a \in A\)）与 B 的不交并为其箭集。\(G_{a}\) 的箭是 a 处的圈；若 \(b \in B\) ，则箭 b 连接 \(\varphi(b)\) 与 \(\psi(b)\) ，而箭 \(c_{1}(b)\) 与 \(c_{2}(b)\) 分别是 \(\varphi(b)\) 与 \(\psi(b)\) 处的圈。箭图 T 将与 \(G_{1}\) 的一个定向树等同起来；它是极大定向树，因为其顶点集等于 \(G_{1}\) 的顶点集（II, 第 157 页, 命题 1）。令 \(a_{0} = a(i_{0})\) 。框架 \(\Gamma_{0}\) （偶 \((\varphi, \psi)\) 的）的箭集既已等同于 B，箭图态射 \(\tau_{1}: \Gamma_{0} \to \text{Grp}(G_{1})\) 便把箭 b 对应到道路类 \(c_{1}(b)^{-1}bc_{2}(b)\) （在图 \(\widetilde{G_{1}}\) 中）。

回顾 \(\theta_{1}\) 表示从 \(\mathrm{G}_{1}\) 到 \(\mathrm{Grp}(\mathrm{G}_{1})\) 的典范箭图态射。以

\[
\lambda \colon \underset {a \in \mathrm{A}} {*} \mathrm{F} (\mathrm{G} _ {a}) * \mathrm{F} (\mathrm{B}) \to \operatorname{Grp} (\mathrm{G} _ {1}) _ {a _ {0}}
\]

表示唯一的群同态，使得

\[
\begin{array}{l l} \lambda (f) = \tau_ {1} (d _ {a}) \theta_ {1} (f) \tau_ {1} (d _ {a}) ^ {- 1} & \text {if } a\in A \text{ 且 } f\in G_{a} \text{;} \\ \lambda (b) = \tau_ {1} (d _ {\varphi (b)}) \tau_ {1} (b) \tau_ {1} (d _ {\psi (b)}) ^ {- 1} & \text {if } b\in B \text{.} \end{array}
\]

于是我们有 \(\lambda = \lambda' \circ \varepsilon\) ，其中 \(\lambda'\) 表示从 \(*_{a \in A} F(G_a) * F(B)\) 到 \(\text{Grp}(G_1)_{a_0}\) 的、由极大定向树 T 所定义的典范群同态（II, 第 179 页, 命题 9），而 \(\varepsilon\) 是群 \(*_{a \in A} F(G_a) * F(B)\) 的唯一自同构，满足 \(\varepsilon(f) = f\) （对 \(a \in A\) 及 \(f \in G_a\) ），以及 \(\varepsilon(b) = c_1(b)^{-1} h_1(b)c_2(b)\) （对 \(b \in B\) ）。由 II, 第 179 页, 注 2，同态 \(\lambda\) 是满射的，其核是 \(*_{a \in A} F(G_a) * F(B)\) 的包含 T 的诸箭的最小正规子群。

以 \(\pi\colon\mathrm{Grp}(\mathrm{G}_{1})\to\mathrm{Coh}(\varphi,\psi)\) 表示典范群胚态射。由 II, 第 177 页, 命题 8 的推论 1，群同态 \(\pi_{a_{0}}\) （从 \(\mathrm{Grp}(\mathrm{G}_{1})_{a_{0}}\) 到 \(\mathrm{Coh}(\varphi,\psi)_{a_{0}}\) ）是满射的，且其核是 \(\mathrm{Grp}(\mathrm{G}_{1})_{a_{0}}\) 的最小正规子群，它包含诸圈 \(\mathrm{Int}(\tau_{1}(d_{a}))( \alpha_{1}(f)\alpha_{1}(g)\alpha_{1}(fg)^{-1})\) （对 \(a\in\mathrm{A}\) 及 \(f,g\in\mathrm{G}_{a}\) ）以及诸圈 \(\mathrm{Int}(\tau_{1}(d_{\varphi(b)}))( \varphi(f)b\psi(f)^{-1}b^{-1})\) （对 \(b\in\mathrm{B}\) 及 \(f\in\mathrm{H}_{b}\) ）。

若以 \(p\colon F(\bigcup G_{a}\cup B)\to \left( \begin{array}{c} *\\ a\in A \end{array}G_{a}\right)*F(B)\) 表示典范满射同态，则我们有 \(\Lambda \circ p = \pi_{a_0}\circ \lambda\) 。这个公式蕴含同态 \(\Lambda\) 是满射的；剩下要确定它的核。

对 \(a \in \mathrm{A}\) 及 \(f \in \mathrm{G}_a\) ，以 {[}f{]} 表示 \(f \in \mathrm{F}(\mathrm{G}_a)\) 在群 \(*_{a \in \mathrm{A}} \mathrm{F}(\mathrm{G}_a) * \mathrm{F}(\mathrm{B})\) 中的像。那么，对 \(a \in \mathrm{A}\) 及 \(f, g \in \mathrm{G}_a\) ，我们有

\[
\operatorname{Int} \left(\tau_ {1} \left(d _ {a}\right)\right) \left(\alpha_ {1} (f) \alpha_ {1} (g) \alpha_ {1} (f g) ^ {- 1}\right) = \lambda ([ f ] [ g ] [ f g ] ^ {- 1}).
\]

类似地，对 \(b \in \mathrm{B}\) 及 \(f \in \mathrm{H}_b\) ，同态 \(\varphi_b\) 与 \(\psi_b\) 的定义（II, 第 193 页, 公式 (6)）蕴含

\[
\begin{array}{r l} & {\mathrm{Int} (\tau_ {1} (d _ {\varphi (b)})) (\varphi (f) h _ {1} (b) \psi (f) ^ {- 1} h _ {1} (b) ^ {- 1})} \\ & {\quad = \tau_ {1} (d _ {\varphi (b)}) \varphi (f) \big (c _ {1} (b) \tau_ {1} (b) c _ {2} (b) ^ {- 1} \big) \psi (f) ^ {- 1}} \\ & {\qquad \qquad \big (c _ {2} (b) \tau_ {1} (b) ^ {- 1} c _ {1} (b) ^ {- 1} \big) \tau_ {1} (d _ {\varphi (b)}) ^ {- 1}} \\ & {\quad = \tau_ {1} (d _ {\varphi (b)}) c _ {1} (b) \varphi_ {b} (f) \tau_ {1} (b) \psi_ {b} (f) ^ {- 1} \tau_ {1} (b) ^ {- 1} c _ {1} (b) ^ {- 1} \tau_ {1} (d _ {\varphi (b)}) ^ {- 1}} \\ & {\quad = \lambda (c _ {1} (b)) \lambda (\varphi_ {b} (f) [ b ] \psi_ {b} (f) ^ {- 1} [ b ] ^ {- 1}) \lambda (c _ {1} (b)) ^ {- 1}.} \end{array}
\]

因此，同态 \(\pi_{a_0} \circ \lambda\) 的核是 \(*\mathrm{F}(\mathrm{G}_a) * \mathrm{F}(\mathrm{B})\) 的最小正规子群，它包含元素 \([f][g][fg]^{-1}\) （对 \(a \in \mathrm{A}\) 及 \(f\) 、\(g \in \mathrm{G}_a\) ）、元素 \(\varphi_b(f)[b]\psi_b(f)^{-1}[b]^{-1}\) （对 \(b \in \mathrm{B}\) 及 \(f \in \mathrm{H}_b\) ）以及元素 \([b]\) （对 \(b \in \mathrm{Fl}(\mathrm{T})\) ）。

最后，同态 \(\Lambda\) 的核是 \((\underset{a\in\mathrm{A}}{*}\mathrm{G}_{a})*\mathrm{F}(\mathrm{B})\) 的最小正规子群，它包含前述诸元素在 \(p\) 下的像，换言之，即元素 \([b]\) （对 \(b\in\mathrm{Fl}(\mathrm{T})\) ）与元素 \(\varphi_{b}(f)[b]\psi_{b}(f)[b]^{-1}\) （对 \(b\in\mathrm{B}\) 及 \(f\in\mathrm{H}_{b}\) ）。命题得证。

